\chapter{अनोखे व्यक्ति}\label{ch:g9:unique-individuals}

सात अरब लोग, और कोई दो एक-से नहीं --- न किसी परिवार में, न किसी स्कूल के
अहाते में, न पूरे इतिहास में; बस समरूप जुड़वाँ आधे अपवाद हैं। तीन अध्याय
इस तथ्य की पूरी मशीनरी जोड़ चुके हैं। यह छोटा-सा समापन अध्याय उस दलील को
एक बार, सिरे से सिरे तक चलाता है और उसके नतीजे बैंक में रख देता है:
\omterm{def:g5:biodiversity:species}{प्रजाति} की विविधता के लिए, ``आनुवंशिक'' सुर्ख़ियाँ पढ़ने के लिए, और इसके
लिए भी कि अनोखा होने का मतलब क्या है और क्या नहीं।

\section{अनोखेपन की उपपत्ति}

\begin{proposition}[कोई दो लोग एक-से क्यों नहीं होते]\label{prop:g9:unique-individuals:theorem}
तीनों कार्य-प्रणालियाँ जोड़ो:
\begin{enumerate}
  \item \emph{आधा करने की फेंटाई}: हर जनक के \omterm{def:g8:sexual-reproduction:gametes}{युग्मक} हर जोड़े से एक
        साथी ढोते हैं, और हर बार अलग-अलग तय करके --- यानी हर जनक से
        $2^{23}$ आधी गड्डियाँ
        (\cref{rem:g9:chromosomes:whyhalve}), और सच पूछो तो उससे कहीं
        ज़्यादा, क्योंकि आधा करने वाले बँटवारे के दौरान जोड़ी के साथी
        आपस में \emph{टुकड़े भी बदल लेते} हैं, और इस तरह गड्डी पूरे
        \omterm{def:g9:chromosomes:chromosomes}{गुणसूत्रों} के दर्जे से भी नीचे कटती है;
  \item \emph{\omterm{def:g5:flower-to-fruit:steps}{निषेचन} की लॉटरी}: उन लाखों में से कौन-सा \omterm{def:g8:reproductive-systems:male}{शुक्राणु}
        (\cref{ex:g8:reproductive-systems:sperm}) किस महीने के \omterm{def:g2:animals-grow-up:born}{अंडे} से
        मिलता है, यह दोनों जगहों को आपस में गुणा कर देता है;
  \item \emph{\omterm{prop:g9:genes-and-dna:explains}{उत्परिवर्तन} की टपकन}: हर पीढ़ी अपनी ओर से कुछ ताज़े
        हिज्जे जोड़ देती है
        (\cref{prop:g9:genes-and-dna:explains})।
\end{enumerate}
मेलों की यह जगह उन तमाम इंसानों की गिनती को बौना कर देती है जो कभी जिए
हैं: हर \omterm{def:g5:flower-to-fruit:steps}{निषेचन} ऐसी गड्डी बाँटता है जो पहले कभी नहीं बाँटी गई, और आगे कभी
नहीं बाँटी जाएगी।
\end{proposition}
\admitted

\begin{example}[जुड़वाँ, और अपवाद की हदें]\label{ex:g9:unique-individuals:twins}
समरूप जुड़वाँ एक ही बँटवारा साझा करते हैं
(\cref{exo:g8:from-egg-to-birth:10}) --- और वे भी अलग होते जाते हैं:
अलग-अलग जीवन अलग-अलग हासिल किए हुए \omterm{def:g9:heredity-and-traits:traits}{लक्षण} लिखते हैं
(\cref{prop:g9:heredity-and-traits:sorting}), \omterm{def:g8:nervous-communication:synapse}{तंत्रिका-संधियों} के
इतिहास अलग होते हैं (\cref{ex:g8:nervous-communication:juggler}), और हर
शरीर की अपनी \omterm{def:g5:first-look-at-cells:cell}{कोशिकाओं} में कुछ अलग देर वाले \omterm{prop:g9:genes-and-dna:explains}{उत्परिवर्तन} भी। अनोखेपन में
पड़ी वह इकलौती दरार ख़ुद ही भर जाती है: बँटवारे के बाद जीना फ़र्क़ पैदा
कर देता है।
\end{example}

\begin{remark}[अनोखे, और फिर भी ज़्यादातर एक-से]\label{rem:g9:unique-individuals:alike}
दोनों सच एक साथ थामो: किन्हीं भी दो इंसानों की \omterm{def:g9:genes-and-dna:dna}{DNA} इबारतें $99$ फ़ीसदी
से ज़्यादा मिलती हैं --- यानी
\cref{prop:g5:first-look-at-cells:principle} वाली एक ही \omterm{def:g5:biodiversity:species}{प्रजाति} की साझी
बनावट --- और तीन अरब अक्षरों में बिखरे हुए वे फ़र्क़ फिर भी सख़्त
अनोखेपन के लिए काफ़ी हैं। ``सब अलग'' और ``सब नातेदार'' एक ही तथ्य के दो
आवर्धन हैं (\cref{met:g9:genes-and-dna:levels}); और पहले सच का दूसरे के
ख़िलाफ़ किया गया हर इस्तेमाल बुरी नैतिकता होने से पहले बुरा जीव विज्ञान
था।
\end{remark}

\section{अनोखापन, काम पर}

\begin{example}[DNA से पहचान]\label{ex:g9:unique-individuals:forensics}
अनोखापन जाँचा जा सकता है, और अदालतें उसे जाँचती हैं: किसी व्यक्ति की
काफ़ी सारी बदलती जगहें, \omterm{def:g5:first-look-at-cells:cell}{कोशिकाओं} के किसी नमूने से पढ़कर, पूरी मानवता में
से एक ही व्यक्ति से मेल खाती हैं (समरूप जुड़वाँ अलग रखे जाएँ तो)। और वही
तुलना परिवार के आवर्धन पर चलाई जाए तो माता-पिता का सवाल हल कर देती है ---
क्योंकि संतान की हर बदलती जगह उसके दोनों जनकों की गड्डियों से बँटी हुई
होनी ही चाहिए। एलबम का तर्क (\cref{ch:g9:heredity-and-traits}) अब एक
औज़ार बन चुका है।
\end{example}

\begin{example}[रक्त चढ़ाना और अंग-प्रत्यारोपण]\label{ex:g9:unique-individuals:medicine}
दवा रोज़ इस अनोखेपन से सौदा करती है। रक्त-समूह
(\cref{ex:g9:genes-and-dna:bloodgroups}) आसान मामला हैं --- चार मोटे
वर्ग, जिन्हें बही से मिलाया जा सकता है। और अंगों का प्रत्यारोपण कठिन
मामला है: प्रतिरक्षा तंत्र
(\cref{ch:g9:immune-defenses} उसकी मशीनरी समझाता है) उससे कहीं महीन
अपनेपन के निशान पढ़ता है, और वे इतने अनोखे होते हैं कि दानी सबसे पहले
परिवार में ढूँढ़े जाते हैं --- \omterm{prop:g6:classification-kinship:kinship}{नातेदारी} जितनी क़रीबी, बँटवारा उतना ही
मिलता-जुलता।
\end{example}

\begin{proposition}[विविधता प्रजाति की बचत है, बैंक में रखी हुई]\label{prop:g9:unique-individuals:reserve}
\Cref{rem:g8:sexual-reproduction:reserve} ने फेंटाई का मक़सद बताने का
वादा किया था; और अब ये कार्य-प्रणालियाँ उसे भुना देती हैं। अनोखे
व्यक्तियों वाली कोई \omterm{def:g5:biodiversity:species}{प्रजाति} अपने लोगों में बिखरी हुई हर मौक़े के लिए
\omterm{def:g9:genes-and-dna:gene}{विकल्पियाँ} थामे रहती है --- किसी बीमारी के ख़िलाफ़ वह रोध जो कुछ ढोते
हैं, और किसी दबाव के ख़िलाफ़ वह सहनशीलता जो कुछ और ढोते हैं
(\cref{ex:g8:sexual-reproduction:bets} वाला बाग़, \omterm{def:g5:biodiversity:species}{प्रजाति} के पैमाने
पर)। और जब \omterm{def:g6:exploring-our-environment:conditions}{परिस्थितियाँ} पलटती हैं, तो
\cref{prop:g8:reproduction-and-environment:limits} वाली रोकें उस
विविधता पर एक-सी नहीं पड़तीं --- और वह असमानता पीढ़ी-दर-पीढ़ी करती क्या
है, ठीक यही अगले दो अध्यायों का विषय है: \omterm{def:g5:biodiversity:species}{प्रजातियों} का बदलना।
\end{proposition}
\admitted

\begin{method}[किसी ``आनुवंशिक'' दावे की जाँच]\label{met:g9:unique-individuals:claims}
\omterm{def:g9:genes-and-dna:gene}{जीनों} और लोगों के बारे में किसी भी सुर्ख़ी के लिए:
\begin{enumerate}
  \item दर्जे की जाँच (\cref{met:g9:genes-and-dna:levels}): यह दावा
        किसी \omterm{def:g9:genes-and-dna:gene}{जीन} के बारे में है, किसी गड्डी के, या किसी परिवार के
        ढर्रे के?
  \item बुनाई की जाँच
        (\cref{prop:g9:heredity-and-traits:sorting}): क्या वह
        विरासत-में-मिली-सीमा-जिसे-जिया-जाता-है वाली बुनाई का लिहाज़
        करता है, या यह दिखावा करता है कि एक \omterm{def:g9:genes-and-dna:gene}{जीन} का मतलब एक नियति है?
  \item विविधता की जाँच: क्या वह किसी \omterm{def:g8:reproduction-and-environment:population}{आबादी} के फैलाव को उस बचत की तरह
        लेता है जो वह है, या किसी ऐसी ``सामान्य'' गड्डी से हटाव की
        तरह, जो होती ही नहीं?
  \item जुड़वाँ की जाँच: क्या वह उन समरूप जुड़वाँ के सामने टिकेगा, जो
        बँटवारा साझा करते हैं और फिर भी अलग होते हैं?
\end{enumerate}
\end{method}

\section{अभ्यास}

\begin{exercise}[$\star$]\label{exo:g9:unique-individuals:1}
अनोखेपन की उपपत्ति की तीनों कार्य-प्रणालियों के नाम बताओ।
\end{exercise}

\begin{exercise}[$\star$]\label{exo:g9:unique-individuals:2}
आधा करने के दौरान टुकड़ों का बदला जाना उस फेंटाई में क्या जोड़ता है?
\end{exercise}

\begin{exercise}[$\star$]\label{exo:g9:unique-individuals:3}
\cref{rem:g9:unique-individuals:alike} के दोनों हिस्से --- वह फ़ीसदी और
वह अनोखापन --- एक ही वाक्य में बताओ।
\end{exercise}

\begin{exercise}[$\star$]\label{exo:g9:unique-individuals:4}
समरूप जुड़वाँ आख़िरकार पहचाने जाने लायक़ अलग कैसे हो जाते हैं?
\end{exercise}

\begin{exercise}[$\star$]\label{exo:g9:unique-individuals:5}
\omterm{def:g9:genes-and-dna:dna}{DNA} की तुलना किसी अदालत के कौन-से दो सवालों का जवाब देती है?
\end{exercise}

\begin{exercise}[$\star$]\label{exo:g9:unique-individuals:6}
प्रत्यारोपण के दानी सबसे पहले परिवार में क्यों ढूँढ़े जाते हैं?
\end{exercise}

\begin{exercise}[$\star\star$]\label{exo:g9:unique-individuals:7}
उस दलील को गुणा करो: दोनों जनकों की आधी गड्डियों वाली जगहें और \omterm{def:g5:flower-to-fruit:steps}{निषेचन} की
लॉटरी जोड़कर ``पहले कभी नहीं बाँटी गई'' वाला नतीजा निकालो।
\end{exercise}

\begin{exercise}[$\star\star$]\label{exo:g9:unique-individuals:8}
\cref{prop:g9:unique-individuals:reserve} को
\cref{exo:g8:sexual-reproduction:10} वाले अंगूर-बाग़ पर भुनाओ: उन
नक़लों के पास वह क्या नहीं था जो \omterm{def:g2:from-seed-to-plant:seed}{बीज} से उगी किसी \omterm{def:g8:reproduction-and-environment:population}{आबादी} के पास होता?
\end{exercise}

\begin{exercise}[$\star\star$]\label{exo:g9:unique-individuals:9}
\cref{met:g9:unique-individuals:claims} इस पर चलाओ: ``स्कूल में
कामयाबी का \omterm{def:g9:genes-and-dna:gene}{जीन} खोज लिया गया''।
\end{exercise}

\begin{exercise}[$\star\star$]\label{exo:g9:unique-individuals:10}
``इंसान की सामान्य गड्डी'' ऐसा वाक्यांश क्यों है जिसका कोई हवाला ही
नहीं? जवाब उस बचत और उस उपपत्ति से दो।
\end{exercise}

\begin{exercise}[$\star\star$]\label{exo:g9:unique-individuals:11}
कोई पितृत्व-जाँच उस शक किए गए पिता को बरी कर देती है जिस पर
\cref{exo:g9:heredity-and-traits:11} वाले पड़ोसी ने ग़लत इल्ज़ाम लगाया
था। समझाओ कि वह जाँच जगह-दर-जगह क्या जाँचती है।
\end{exercise}

\begin{exercise}[$\star\star\star$]\label{exo:g9:unique-individuals:12}
``सब अलग, सब नातेदार, और दोनों एक ही कार्य-प्रणाली से।'' यह समापन
अनुच्छेद लिखो: फेंटाई, \omterm{prop:g9:genes-and-dna:explains}{उत्परिवर्तन} और साझा कूट, हर एक अपनी जगह पर।
\end{exercise}

\section{समस्या: पुराने मुक़दमे की संगोष्ठी}

\begin{problem}[{सप्ताहांत समस्या --- एक बाल, तीन सवाल, तीस
साल}]\label{pb:g9:unique-individuals:1}
एक (काल्पनिक) पुराने मुक़दमे की संगोष्ठी: तीस साल पुरानी एक अनसुलझी
चोरी, \omterm{def:g1:plants-around-us:parts}{जड़} की \omterm{def:g5:first-look-at-cells:cell}{कोशिकाओं} समेत सँभालकर रखा एक बाल, और तुलना के आधुनिक
तरीक़े। यह संगोष्ठी विद्यार्थियों को दिखाती है कि \omterm{def:g9:genes-and-dna:dna}{DNA} से पहचान क्या कह
सकती है और क्या नहीं।

\medskip\noindent\textbf{भाग I --- वह बाल थामे क्या है।}
\begin{enumerate}
  \item पूरी पढ़त के लिए उस बाल की \emph{\omterm{def:g1:plants-around-us:parts}{जड़}} का सँभला होना क्यों
        ज़रूरी है (\cref{def:g6:cell-unit-of-life:parts} जानता है कि
        \omterm{def:g9:genes-and-dna:dna}{DNA} रहता कहाँ है)?
  \item प्रयोगशाला बहुत ज़्यादा बदलती जगहों का एक सेट पढ़ती है।
        बदलती जगहें ही क्यों --- वे 99 फ़ीसदी साझी जगहें क्या साबित
        करतीं?
  \item वह ब्योरा हर जगह पर एक संदिग्ध से मेल खाता है। बताओ कि यह मेल
        क्या दावा करता है, और उसका इकलौता व्यवस्थित अपवाद भी
        (\cref{ex:g9:unique-individuals:twins})।
  \item बचाव पक्ष बताता है कि उस संदिग्ध का एक समरूप जुड़वाँ विदेश में
        है। अब अदालत को क्या करना होगा, और क्यों?
\end{enumerate}

\medskip\noindent\textbf{भाग II --- परिवार वाले सवाल।}
\begin{enumerate}[resume]
  \item पहले लिया गया एक अधूरा ब्योरा ``क वंश के किसी क़रीबी नातेदार''
        की ओर इशारा करता था। समझाओ कि कोई \emph{अधूरा} मेल \omterm{prop:g6:classification-kinship:kinship}{नातेदारी}
        कैसे बताता है
        (\cref{prop:g9:unique-individuals:theorem} वाला बँटवारा,
        उलटा चलाया हुआ)।
  \item एक जूरी-सदस्य पूछता है कि क्या वह ब्योरा संदिग्ध के \omterm{def:g9:heredity-and-traits:traits}{लक्षण}
        बताता है --- शक्ल, मिज़ाज, ``अपराध वाले \omterm{def:g9:genes-and-dna:gene}{जीन}''। इस सवाल पर
        \cref{met:g9:unique-individuals:claims} चलाओ, जाँच-दर-जाँच।
  \item एक और जूरी-सदस्य पूछता है कि क्या वह मेल ``इस शहर में एक बार''
        भर हो सकता है, न कि पूरी मानवता में एक बार। पढ़ी जाने वाली
        जगहों की गिनती तय क्या करता है?
  \item संगोष्ठी का संचालक चेताता है: ``जीव विज्ञान बताता है कि
        \omterm{def:g5:first-look-at-cells:cell}{कोशिकाएँ} किसकी हैं; और यह मुक़दमे को अब भी बताना है कि वे वहाँ
        पहुँचीं कैसे।'' दोनों दावों में फ़र्क़ करो।
\end{enumerate}

\medskip\noindent\textbf{भाग III --- संगोष्ठी का समापन।}
\begin{enumerate}[resume]
  \item उस औज़ार का सार एक वाक्य में दो: किन दो अध्यायों की
        कार्य-प्रणालियाँ (\cref{ch:g9:chromosomes},
        \cref{ch:g9:genes-and-dna}) किसी ब्योरे को अनोखा और किसी
        परिवार को पढ़ने लायक़ बनाती हैं।
  \item वही औज़ार तीस साल पुराने बाल पर भी काम क्यों करता है? (\omterm{def:g9:genes-and-dna:dna}{DNA} का
        कौन-सा गुण --- \cref{prop:g9:genes-and-dna:explains} वाली
        नक़ल को छोड़कर --- यहाँ सहारा बना हुआ है: टिकाऊपन।)
  \item संगोष्ठी नैतिकता पर ख़त्म होती है: अनोखेपन के वे दो इस्तेमाल
        गिनाओ जिन्हें यह अध्याय लोगों की सेवा करते दिखाता है, और वह
        ऐतिहासिक दुरुपयोग भी जिससे
        \cref{rem:g9:unique-individuals:alike} चेताती है।
  \item समापन उस उपपत्ति को अदालत के लिए दोबारा कहकर करो: वह क्या है
        जो कभी दो बार नहीं हुआ, और क्यों।
\end{enumerate}
\end{problem}
